\section{Experimental setup}
\label{sec:setup}

\subsection{Testbeds}

\paragraph{Two-scale Lorenz-96.}
Eight slow variables coupled to $J$ fast variables each, integrated with RK4.
The assimilation window is 500 steps. Observations cover a 24-dimensional
subspace; the full physical state is 40-dimensional. Every window draws its own
physical parameters, randomised by $\pm 20\%$, so that no method can memorise a
single parameter setting.

\paragraph{Two-layer quasi-geostrophic channel.}
A Phillips-type two-layer channel, integrated in potential vorticity $q$ with a
spectral inversion for the streamfunction $\psi$, forced by a moving-storm wind
stress curl. Observations are sparse meridional columns, each observed once per
day at its own randomly sampled intra-day step.

\subsection{The model-error axis: \Sz{} and \So{}}
\label{sec:s0s1}

Two scenarios are used throughout:
\begin{description}[leftmargin=2em,style=nextline]
  \item[\Sz{} (perfect model).] The assimilating model's parameters equal those
    used to generate the truth for that window.
  \item[\So{} (model error).] The assimilating model's parameters carry an
    additional $\pm 10\%$ bias. The truth is drawn from the same distribution in
    both scenarios; only the \emph{model handed to the estimator} changes.
\end{description}

\paragraph{A caveat that must be stated rather than cashed in.}
The \Sz{}/\So{} axis prices model error \emph{only for methods that run a forward
model at inference}. A learned estimator that receives no parameters and
integrates nothing is definitionally unaffected, and its $\So/\Sz \approx 1.00$
is not a robustness result. We therefore separate three method classes
throughout (Table~\ref{tab:classes}) and never table a model-free ratio as
cross-class robustness evidence.

\begin{table}[t]
\centering
\small
\begin{tabular}{@{}p{0.22\textwidth}p{0.36\textwidth}p{0.34\textwidth}@{}}
\toprule
\textbf{class} & \textbf{what \So{} changes for it} & \textbf{a ratio $\approx 1.00$ means} \\
\midrule
classical DA (4D-Var, EnKF/ETKF) & integrates a biased forward model; the bias
compounds through the integration & n/a --- they degrade by $1.68$--$1.80\times$ \\
model-free learned & nothing: no parameters, no forward model & \emph{definitional};
carries no information \\
conditioned learned & the conditioning inputs are biased or corrupted & a genuine
robustness result \\
\bottomrule
\end{tabular}
\caption{What \So{} means depends on whether, and how, a method consumes model
information. This distinction governs every cross-class statement below.}
\label{tab:classes}
\end{table}

\subsection{Metrics}
RMSE is reported pooled over all windows and timesteps
($\sqrt{\overline{(\cdot)^2}}$), the same convention for every method. Explained
variance (EV) is reported per field. For ensemble methods the spread/skill ratio
is compared against the finite-ensemble reliability target
$\sqrt{(N-1)/(N+1)}$, and proper ensemble CRPS / energy scores are computed at a
single uniform ensemble-size convention ($N=30$).

\needsrun{Rank histograms are the primary posterior-diagnostic figure for a DA
readership and are not yet implemented anywhere in the codebase. This is the
largest new-code item remaining (experiment D5).}
